\section*{Hoofdstuk \ref{ch:g12:work-and-energy} --- Arbeid en mechanische energie}

\begin{solution}{exo:g12:work-and-energy:1}
$E_p = \tfrac12 \times 200 \times 0.050^2 = \qty{0.25}{J}$. Bij dubbele
uitrekking: $\times 4$, dus \qty{1.0}{J}. De tweede \qty{5.0}{cm} alleen kost
$1.0 - 0.25 = \qty{0.75}{J}$ --- daar is de \omterm{def:g10:inertia:force}{kracht} groter.
\end{solution}

\begin{solution}{exo:g12:work-and-energy:2}
$W = mg(z_A - z_B) = 0.90 \times 9.81 \times (12 - 30) \approx
\qty{-1.6e2}{J}$, wat de route ook is --- dus hetzelfde langs de rechte klim.
Heen en terug: \qty{0}{J} (een \omterm{def:g12:work-and-energy:conservative}{conservatieve kracht}).
\end{solution}

\begin{solution}{exo:g12:work-and-energy:3}
(a), (b) conservatief: de \omterm{def:g11:work-of-force:work}{arbeid} ligt vast door de uiteinden (hoogteverschil;
uitrekking). (c), (e) niet-conservatief: $-fL$ en de luchtweerstand rekenen per
afgelegde \omterm{def:g10:orders-of-magnitude:unit}{meter}. (d) arbeidsloos: altijd loodrecht op de beweging.
\end{solution}

\begin{solution}{exo:g12:work-and-energy:4}
$h = L(1 - \cos\qty{60}{\degree}) = \qty{1.0}{m}$;
$v = \sqrt{2 \times 9.81 \times 1.0} \approx \qty{4.4}{m/s}$. De \omterm{def:g10:inertia:inventory}{spankracht}
staat op elk ogenblik loodrecht op de beweging: nul \omterm{def:g11:work-of-force:power}{vermogen}, nul \omterm{def:g11:work-of-force:work}{arbeid}.
\end{solution}

\begin{solution}{exo:g12:work-and-energy:5}
$F = P/v = 250/9.0 \approx \qty{28}{N}$. Bij \qty{5.0}{m/s}:
$F = 250/5.0 = \qty{50}{N}$ (grotendeels de component van de zwaartekracht
langs de helling).
\end{solution}

\begin{solution}{exo:g12:work-and-energy:6}
$E_p = \tfrac12 \times 450 \times 0.080^2 = \qty{1.44}{J}$;
$v = \sqrt{2 \times 1.44/0.020} = \qty{12}{m/s}$;
$h = E_p/mg = 1.44/(0.020 \times 9.81) \approx \qty{7.3}{m}$.
\end{solution}

\begin{solution}{exo:g12:work-and-energy:7}
Trapezia (stappen van \qty{2.0}{cm}): $(0.03 + 0.10 + 0.19 + 0.30) =
\qty{0.62}{J}$. Model met een rechte: $k = 18.0/0.080 = \qty{225}{N/m}$,
$\tfrac12 kx^2 = \qty{0.72}{J}$. De gemeten kromme zakt bij kleine $x$ onder
de rechte, dus is de ware oppervlakte kleiner.
\end{solution}

\begin{solution}{exo:g12:work-and-energy:8}
$\Delta E_m = \tfrac12 \times 55 \times 9.0^2 - 55 \times 9.81 \times
8.0 = 2228 - 4316 \approx \qty{-2.1e3}{J}$;
$f = 2089/25 \approx \qty{84}{N}$. Het \omterm{def:g10:universal-gravitation:weight}{gewicht} is conservatief en de
\omterm{def:g12:newtons-laws:friction}{normaalkracht} arbeidsloos: alleen de daling en de \emph{lengte} van de weg
tellen mee, niet de vorm.
\end{solution}

\begin{solution}{exo:g12:work-and-energy:9}
$H_{\min} = \tfrac52 R = \qty{0.50}{m}$. Van \qty{60}{cm}:
$v_{\text{top}} = \sqrt{2 \times 9.81 \times (0.60 - 0.40)} \approx
\qty{2.0}{m/s}$, en $v_{\text{top}}^2/(gR) = 2.0$ --- het dubbele van het
minimum. De \omterm{def:g10:inertia:inventory}{wrijving} schept onderweg \omterm{def:g10:energy-conservation:energy}{energie} af, dus hebben echte banen extra
hoogte nodig.
\end{solution}

\begin{solution}{exo:g12:work-and-energy:10}
\omterm{def:g12:work-and-energy:diagram}{Keerpunten}: $8.0\,x^2 = 2.0$, dus $x = \pm\qty{0.50}{m}$. Maximale snelheid in
$x = 0$: $v = \sqrt{2 \times 2.0/0.100} \approx \qty{6.3}{m/s}$. Een heen en
weer gaande trilling in een parabolische put --- de $E_p$ van de \omterm{prop:g12:oscillators-and-time:spring-period}{veeroscillator}
heeft precies deze vorm.
\end{solution}

\begin{solution}{exo:g12:work-and-energy:11}
$\tfrac12 kx^2 = \tfrac12 mv^2$:
$x = v\sqrt{m/k} = 2.0\sqrt{1200/\num{6.0e5}} \approx \qty{8.9}{cm}$. Bij
\qty{4.0}{m/s}: \qty{18}{cm} --- het dubbele, niet het viervoud: de \omterm{def:g10:energy-conservation:energy}{energie}
verviervoudigt, maar $E_p$ groeit als $x^2$, dus verdubbelt $x$ slechts.
\end{solution}

\begin{solution}{exo:g12:work-and-energy:12}
(a) $E_m = \qty{5.0}{J} < \qty{6.0}{J}$: het raakt er niet over; het stopt op
de flank waar $E_p = \qty{5.0}{J}$ en rolt terug. (b) Heuveltop:
$E_k = \qty{1.0}{J}$, $v = \sqrt{2 \times 1.0/0.50} = \qty{2.0}{m/s}$; dal:
$E_k = \qty{5.0}{J}$, $v \approx \qty{4.5}{m/s}$. (c) $x = 0$: een \omterm{prop:g12:work-and-energy:equilibria}{labiel
evenwicht} (maximum); $x = \qty{1.5}{m}$: een \omterm{prop:g12:work-and-energy:equilibria}{stabiel evenwicht} (minimum).
\end{solution}

\begin{solution}{exo:g12:work-and-energy:13}
$\tfrac12 mv^2 = mgR_E$:
$v = \sqrt{2 \times 9.81 \times \num{6.37e6}} \approx \qty{1.12e4}{m/s}
= \qty{11.2}{km/s}$. Elke term is evenredig met $m$: sonde en steentje gelijk.
Net eronder klimt de lancering nog enorm ver, komt een \omterm{def:g12:work-and-energy:diagram}{keerpunt} tegen en valt
terug: gebonden.
\end{solution}

\begin{solution}{exo:g12:work-and-energy:14}
$E_m = \tfrac12 kA^2 = \tfrac12 \times 40 \times 0.060^2 =
\qty{7.2e-2}{J}$; $v_{\max} = \sqrt{2E_m/m} =
\sqrt{2 \times 0.072/0.250} \approx \qty{0.76}{m/s}$; $E_k = E_p$ waar
$\tfrac12 kx^2 = E_m/2$: $x = \pm A/\sqrt 2 \approx \pm\qty{4.2}{cm}$.
\end{solution}

\begin{solution}{exo:g12:work-and-energy:15}
$h = v^2/2g = 9.5^2/19.62 \approx \qty{4.6}{m}$; de lat op $\approx
4.6 + 1.0 = \qty{5.6}{m}$. Het record ligt hoger: springers voegen met armen
en benen \omterm{def:g11:work-of-force:work}{arbeid} toe op de buigende stok. Maar de energieschatting zet de
schaal --- geen enkele sprint brengt een polsstokspringer op \qty{10}{m}.
\end{solution}

\begin{solution}{pb:g12:work-and-energy:1}
\textbf{1.} $v = \sqrt{2gH} = \sqrt{2 \times 9.81 \times 45} \approx
\qty{29.7}{m/s}$: het \omterm{def:g10:universal-gravitation:weight}{gewicht} is conservatief en de \omterm{def:g12:newtons-laws:friction}{normaalkracht}
arbeidsloos --- het profiel valt weg.

\textbf{2.} Poort: $E_m = mgH = 70 \times 9.81 \times 45 \approx
\qty{3.09e4}{J}$. Tafel: $E_m = \tfrac12 \times 70 \times 26^2
\approx \qty{2.37e4}{J}$.

\textbf{3.} $\Delta E_m \approx \qty{-7.2e3}{J}$; de
mechanische-energiestelling schrijft dat op het conto van de
\omterm{def:g12:work-and-energy:conservative}{niet-conservatieve krachten}: de \omterm{def:g10:inertia:inventory}{wrijving} met de sneeuw en de luchtweerstand.

\textbf{4.} $f = \num{7.2e3}/90 \approx \qty{80}{N}$.

\textbf{5.} $\num{7.2e3}/\num{3.09e4} \approx 23\%$. Wax verlaagt $f$,
ineengedoken zitten verlaagt de luchtweerstand: beide verkleinen $fL$, de enige
onderhandelbare term.

\textbf{6.} \omterm{def:g12:projectile-motion:freefall}{Vrije val} met een horizontale beginsnelheid --- een worp.
$v_x = \qty{26}{m/s}$; $v_y = \sqrt{2 \times 9.81 \times 40} \approx
\qty{28.0}{m/s}$.

\textbf{7.} $v = \sqrt{26^2 + 28.0^2} = \sqrt{1461} \approx
\qty{38.2}{m/s}$; met \omterm{def:g10:energy-conservation:energy}{energie}: $v = \sqrt{v_0^2 + 2gh} = \sqrt{676 + 785}$
--- identiek. Ongeveer \qty{138}{km/h}.

\textbf{8.} $t = v_y/g = 28.0/9.81 \approx \qty{2.9}{s}$;
$d = v_0 t \approx \qty{74}{m}$.

\textbf{9.} Nee. De liftkracht staat loodrecht op de snelheid (arbeidsloos) en
de luchtweerstand werkt haar tegen (negatieve \omterm{def:g11:work-of-force:work}{arbeid}): $\Delta E_m \leq 0$, dus kan
de echte landingssnelheid alleen \emph{kleiner} zijn --- de lucht rekt de
vlucht op, ze voedt haar niet.

\textbf{10.} $\alpha = \arctan(28.0/26) \approx \qty{47}{\degree}$ onder de
horizontaal.

\textbf{11.} Hoek met de helling: $47 - 35 = \qty{12}{\degree}$.
$v_\parallel = 38.2\cos\qty{12}{\degree} \approx \qty{37.4}{m/s}$;
$v_\perp = 38.2\sin\qty{12}{\degree} \approx \qty{7.9}{m/s}$.

\textbf{12.} Opgevangen: $\tfrac12 \times 70 \times 7.9^2 \approx
\qty{2.2e3}{J}$, tegenover $\tfrac12 \times 70 \times 38.2^2 \approx
\qty{5.1e4}{J}$ in totaal: ongeveer $4\%$.

\textbf{13.} $h_{\text{eq}} = 7.9^2/19.62 \approx \qty{3.2}{m}$ --- een
stoutmoedige sprong van een muurtje. Vlak: $38.2^2/19.62 \approx \qty{74}{m}$
--- niet te overleven.

\textbf{14.} Een helling evenwijdig aan de vluchtbaan houdt $v_\perp$ klein,
zodat de benen \omterm{def:g10:orders-of-magnitude:unit}{meters} opvangen en niet de hele berg \omterm{def:g11:mechanical-energy:kinetic}{kinetische energie}.

\textbf{15.} $H'_0 = v_0'^2/2g = 576/19.62 \approx \qty{29.4}{m}$.

\textbf{16.} $mgH' - f\,\dfrac{H'}{\sin\qty{35}{\degree}} =
\tfrac12 m v_0'^2$.

\textbf{17.} $H' = \dfrac{\tfrac12 \times 70 \times 576}
{686.7 - 80/0.574} = \dfrac{\num{20160}}{547} \approx \qty{37}{m}$.

\textbf{18.} $36.8/29.4 \approx 1.25$: een toeslag van $25\%$.

\textbf{19.} Voor \qty{60}{kg}: $H' = \num{17280}/(588.6 - 139.5)
\approx \qty{38.5}{m}$. De \emph{lichtere} springer heeft de hoogste toren
nodig: de wrijvingsrekening $fL$ is dezelfde, maar haar energiebudget $mgH'$
is kleiner.

\textbf{20.} Kondig een toren van \qty{37}{m} aan voor \qty{24}{m/s} op de
tafel --- en beken dat een kwart ervan, zo'n \qty{7}{m} beton, de toeslag van
de \omterm{def:g10:inertia:inventory}{wrijving} op de wrijvingsloze droom is.
\end{solution}
